\scanpage{16}
\mspage{\handwritten{16\unsure{a|.}}}
algébr\handwritten{iq}uement clos. Mais alors, $s$ désignant le point fermé de $S$, comme $G$ est étale et séparé sur $S$, il contient un sous-schéma à la fois ouvert et fermé $G'$, \inserted{fini sur $S$,} tel que $G'_s=G_s$ \inserted{(EGA II 6.2.6.)}, et on voit de suite que $G'$ est ici un sous-groupe de $G$. \deleted{\unsure{Par suite|Ensuite} la fibre géométrique $f(s)$} D'ailleurs $G'$ \deleted{\unsure{est|qui}} étant étale fini sur $S$ \deleted{\unsure{qui}} $=\operatorname{Spec}(V)$, $V$ complet à corps résiduel algébriquement clos, est un groupe constant, donc de la forme $A_S$, où $A\handwritten{=}G'(k(s))=G(k(s))$ a pour classe $f(s)$. \deleted{On voit donc} Si $B$ est la fibre géométrique de $G$ en le point générique $t$ de $S$, on a donc un \deleted{\unsure{isomor}}\inserted{monomor}phisme canonique $A\handwritten{\to}B$, ce qui prouve $f(s)\handwritten{\leq} f(t)$. (NB cette démonstration prouve en fait la semi-continuité pour une relation d'ordre sur $E$ plus fine que celle indiquée dans \handwritten{2}.1.). Le fait que $f$ soit continue en $s$ si et seulement si $s\handwritten{\rightsquigarrow}\operatorname{card} f(s)$ l'est résulte du fait que pour $w,w'\handwritten{\in}E$, les relations \deleted{\unsure{$w$|$w,$}} $w\handwritten{\leq}w'$ et $\operatorname{card} w=\operatorname{card} w'$ impliquent $w=w'$. Le fait que cette condition soit équivalente à\deleted{\unsure{s|,}} la finitude de $W$ sur un voisinage de $s$ est alors indépendant de\deleted{\unsure{s}} \inserted{la} structure\deleted{\unsure{s}} de groupe\deleted{\unsure{s}} \inserted{\handwritten{sur $W$}}, et a été signalé après \inserted{\handwritten{Exp $\overline{\underline{\mathrm{XI}}}$}} 5.10. ; d'ailleurs sa démonstration \deleted{\unsure{par}} se fait aisément par les arguments précédents, \inserted{en} utilisant le c\handwritten{ri}tère valuatif de propreté EGA II 7.3.8.

\uline{Lemme} \handwritten{2}.3. Soient $G$ un groupe algébrique affine lisse sur un corps algébriquement clos $k$, $R\handwritten{\subset}T$ deux sous-tores, $W(R)$ et $W(T)$ les deux groupes finis associés comme dans \inserted{\handwritten{Exp $\overline{\underline{\mathrm{XI}}}$,}} 5.9. , quotient\handwritten{s} du normalisateu\unsure{r} par le centralisateur. Alors $W(R)$ est isomorphe à un quotient d'un sous-groupe de $W(T)$.

\deleted{Dans la démonstration qui suit, il sera commode d'interpréter $G,R,T,C(R),N(R),C(T),N(T)$ de façon purement ensembliste, comme les groupes de points des préschémas $G,R,T$ dans $k$, et les centralisateurs et nor}
% archive sheet number: 15
